% Generated draft: requires mathematical and source review.
\documentclass[11pt]{article}
\usepackage[margin=25mm]{geometry}
\usepackage{fontspec}
\setmainfont{Latin Modern Roman}
\usepackage{amsmath,amssymb,mathtools,unicode-math}
\setmathfont{Latin Modern Math}
\usepackage{enumitem,needspace,xurl,hyperref}
\setlength{\parindent}{0pt}
\setlength{\parskip}{5pt}
\setlength{\emergencystretch}{3em}
\newcommand{\N}{\mathbb{N}}
\newcommand{\Z}{\mathbb{Z}}
\newcommand{\Q}{\mathbb{Q}}
\newcommand{\R}{\mathbb{R}}
\newcommand{\C}{\mathbb{C}}
\newcommand{\F}{\mathbb{F}}
\newcommand{\A}{\mathbb{A}}
\newcommand{\PP}{\mathbb{P}}
\newcommand{\E}{\mathbb{E}}
\newcommand{\eps}{\varepsilon}
\newcommand{\dd}{\,\mathrm{d}}
\begin{document}
\section*{UnsolvedMath ID: 6200013}
\textbf{Source problem code:} catalogue entry
\subsection{Definitions and mathematical statement}
An ideal boundary records asymptotic directions in a space with a specified geometry. For a proper geodesic hyperbolic space, its Gromov boundary can be represented by geodesic rays modulo finite Hausdorff distance, with the boundary topology. For a word-hyperbolic group one uses a Cayley graph. Other spaces, for example CAT(0) spaces, may instead use a visual boundary whose conventions must be fixed. The source does not restrict all its examples to one boundary construction.

\textbf{Research problem (Tadeusz Januszkiewicz).} Find further universality phenomena: classes of groups or spaces of different nature whose ideal boundaries are nevertheless all homeomorphic, beyond the examples involving Menger compacta, right-angled hyperbolic buildings, and Davis--Vinberg complexes.

A precise form of a proposed example would give a class $\mathcal C$, define $\partial_\infty X$ for its members, identify a compactum $B$, and prove
\[
 \partial_\infty X\cong B\qquad\text{for every }X\in\mathcal C.
\]
The source seeks genuinely further examples but supplies no criterion for ``different nature,'' such as nonisomorphism or non-quasi-isometry, and no exhaustive target class. Neither criterion is added. Universality here means a common boundary homeomorphism type across a class; it is not automatically the separate assertion that every compactum embeds into that boundary.

The supplied summary mentions right-angled Coxeter examples with Menger or spherical boundaries and Menger-curve boundaries for certain random-group models, but gives no hypotheses or citations. Its phrase about ``dimension-free compacta'' does not specify a valid dimensional embedding assertion and is not imposed as part of this problem. These reports provide tentative examples/progress rather than a complete classification or a precise remaining conjecture.
\subsection{Short English statement}
Find new, structurally different families of groups or spaces sharing one ideal-boundary homeomorphism type, beyond the listed Menger/building/Davis--Vinberg examples. The boundary construction, diversity criterion, and scope are open-ended; common boundary type is distinguished from universal embedding properties.
\subsection{Sources}
\begin{itemize}
\item Local UnsolvedMath registry, permanent ID 6200013, source code catalogue entry, ``Boundaries of Groups and Kleinian Groups --- Problem 13''; proposer: Tadeusz Januszkiewicz, proposed year 2005. Status partially solved, classification PARTIAL-PROGRESS, updated 2026-08-06.
\item Kapovich, \emph{Problems on Boundaries of Groups and Kleinian Groups} (2005), Problem 13, PDF page 5, \url{https://www.math.ucdavis.edu/~kapovich/EPR/problems.pdf}. Reported PDF-text-normalized extraction and access: 2026-07-29. The source labels it a problem; current status NEEDS\_REVIEW. Rights: public author-hosted PDF; statement-reuse terms NEEDS\_REVIEW.
\item The statement names Menger compacta, right-angled hyperbolic buildings, and Davis--Vinberg complexes. The summary additionally mentions right-angled Coxeter and random-group examples without primary references or exact assumptions.
\end{itemize}
\end{document}
